\section{Memory Abstraction}
\subsection{Memory Management}
\paragraph{Random Access Memory (RAM)}
\begin{itemize}
  \item \textcolor{Blue}{Volatile} memory: data is lost when power is turned off
  \item \textcolor{Blue}{Byte-addressable}: each byte has a unique address
  \item \textcolor{Blue}{Directly accessible} by CPU: CPU can read/write to RAM without
    any intermediate steps
\item \textcolor{Blue}{\emph{Contiguous} memory region}: consists an
    interval of \textcolor{Bittersweet}{consecutive} addresses
\end{itemize}
In a process, there are generally two types of data:
\begin{itemize}
  \item \textcolor{Blue}{Transient Data}: Data that is valid only for a limited duration
  \item \textcolor{Blue}{Persistent Data}: Data that is valid for the entire
    duration of the process unless explicitly deallocated
\end{itemize}
The OS \textcolor{Bittersweet}{manages the memory space} of processes:
\begin{itemize}
  \item Allocate memory space for new processes
  \item Manage memory space for processes
  \item Protect memory space of processes from each other
  \item Provides memory related system calls
  \item Manage memory space for internal OS data structures
\end{itemize}

\textcolor{ForestGreen}{\textbf{Without memory abstraction}}, physical
addresses are directly used. Memory access is straightforward, but
\begin{itemize}
  \item Processes might not have \textcolor{Bittersweet}{contiguous memory
    space}
  \item Memory addresses may not be \textcolor{Bittersweet}{consistent across
    runs} (and do not begin at 0)
  \item \textcolor{Bittersweet}{Hard to protect memory space} of processes from
    each other, as they can directly access physical addresses of other
    processes
\end{itemize}
\textbf{\textcolor{Blue}{Address Relocation}}: Map all memory references of a
process to the actual physical addresses. \textcolor{ForestGreen}{e.g. Add a
constant offset} to all memory references. Results in
\textcolor{Bittersweet}{slow loading time}, and memory references are
\textcolor{Bittersweet}{difficult to distinguish} from integer constants (in
binary form).

\textbf{\textcolor{Blue}{Base and Limit Registers}}: Special registers in
\textcolor{Blue}{memory context} storing the offset of all memory references
(starting address) and the range (length of memory space). Incurs an
\textcolor{Bittersweet}{additional addition and comparison} for each memory
access (even for instruction retrieval).

\paragraph{Logical Address}
Having a \textbf{\textcolor{Blue}{logical interface}} for memory accesses to
processes, where each process has its own \textcolor{Bittersweet}{logical
address space that is mapped to the physical address space} by the OS.

\subsection{Contiguous Memory Management}
Assumptions:
\begin{itemize}
  \item Processes (and data) must be loaded into memory for execution
    (\textcolor{Bittersweet}{Store Memory concept})
  \item Instructions operate solely on registers except for load and store that
    access memory (\textcolor{Bittersweet}{Load-Store Memory})
  \item Processes require \textbf{contiguous memory space} (Memory Paritition)
  \item Physical memory is \textcolor{Bittersweet}{large enough} to hold all processes
    \begin{itemize}
      \item Removing terminated processes
      \item Swapping out blocked processses to secondary storage
    \end{itemize}
\end{itemize}
\paragraph{Partition allocation schemes}
\textcolor{Blue}{\textbf{Fixed size} partitioning}: Memory is divided into
fixed number of partitions. Each partition can hold one process.
\vspace{-4mm}
\begin{multicols}{2}
  \begin{itemize}
    \item \textcolor{ForestGreen}{Easy to manage}
    \item \textcolor{ForestGreen}{Fast} to allocate and deallocate
  \end{itemize}
  \vfill\null
  \begin{itemize}
    \item Partition size must be \textcolor{Bittersweet}{large enough} to hold
      the largest process
    \item Can lead to \textcolor{Bittersweet}{internal fragmentation} (wasted
      space within partitions)
  \end{itemize}
\end{multicols}

\textcolor{Blue}{\textbf{Variable size} partitioning}: Memory is allocated
dynamically based on the size of the process.
\vspace{-4mm}
\begin{multicols}{2}
  \begin{itemize}
    \item More \textcolor{ForestGreen}{efficient} use of memory
    \item More \textcolor{ForestGreen}{flexible}, able to accommodate processes
      of varying sizes
    \item Eliminates \textcolor{ForestGreen}{internal fragmentation}
  \end{itemize}
  \vfill\null
  \begin{itemize}
    \item Can lead to \textcolor{Bittersweet}{external fragmentation} (large
      number of ``holes'', each too small for new processes)
    \item \textcolor{Bittersweet}{More information} to manage in OS
    \item \textcolor{Bittersweet}{Slower} to allocate and deallocate
  \end{itemize}
\end{multicols}

\vspace{-4mm}
\subsection{Dynamic Partitioning}
\paragraph{Allocation}
A list of partitions and holes is maintained. When a partition of size $N$ is
needed, OS searches for a suitable hole of size $M>N$:
\begin{itemize}
  \item \textcolor{Blue}{First-Fit}: Allocate the first hole that is large enough.
  \item \textcolor{Blue}{Best-Fit}: Allocate the smallest hole that is large enough.
  \item \textcolor{Blue}{Worst-Fit}: Allocate the largest hole.
\end{itemize}
and splits the hole into $N$ and $M-N$.
\paragraph{Deallocation}
When a process terminates, the occupied partition is freed and becomes a hole.
OS can reduce fragmentation by:

\textcolor{Blue}{\textbf{Merging}}: Combine adjacent free partitions into a
single larger hole.

\textcolor{Blue}{\textbf{Compaction}}: Move partitions to consolidate holes and
eliminate fragmentation. \textcolor{Bittersweet}{Expensive operation}, cannot
be invoked too frequently.
